\chapter{Implementation and Validation}
\section{Introduction}
This chapter will provide the validation of the architecture designed in Chapter 3 for the BGP EVPN-VXLAN IP Fabric solution. Each component of the solution is validated using a unique test scenario based on a methodology comprising six elements – objectives, topology context, configuration, expected outcome, findings, and conclusion. These eight tests are done in a methodological sequence starting from underlay convergence all the way to resilience testing, validating that all issues highlighted in Chapter 1 have been fixed by this architecture.

\section{Implementation \& Validation}
\subsection{Installation and global Configuration for fabric nodes}
\begin{figure}[!htbp]
\centering
\begin{lstlisting}[
    language=bash,
    style=nexus,
]
feature ospf
feature bgp
feature bfd
feature interface-vlan
feature vn-segment-vlan-based
feature nv overlay
nv overlay evpn

bfd interval 300 min_rx 300 multiplier 3
fabric forwarding anycast-gateway-mac 0200.5e00.01ff
\end{lstlisting}
\caption{Global feature configuration for underlay and overlay}
\label{fig:global-config}
\end{figure}

\subsection{Scenario 1: Underlay Convergence and ECMP Verification}
\paragraph{Objective:}~\\
Confirm the establishment of the OSPF adjacency among all the Spine and Leaf nodes at both sites, reachability of all the loopbacks through the IGP, and load balancing across both Spines at the same time using ECMP, verifying that there are no blocked or stand-by links and thus that all the fabric’s bandwidth is used for packet forwarding.

\paragraph{Topology Context:}~\\
In this configuration, all the Spine and Leaf nodes at both sites are present. All the Leaf nodes are connected to both Spine nodes using /31 OSPF interfaces within Area 0. BFD is enabled on all OSPF neighborships with a detection time of 300 milliseconds and a multiplier of 3.

\paragraph{Configuration Applied:}~\\
The following key commands represent the core underlay configuration applied to each device type.
\begin{figure}[!htbp]
\centering
\begin{lstlisting}[
    language=bash,
    style=nexus,
]
interface loopback0
  ip address 10.1.1.1/32
  ip router ospf 1 area 0.0.0.0

interface loopback1
  ip address 10.2.1.1/32
  ip router ospf 1 area 0.0.0.0

interface Ethernet1/1
  description to-Spine1-1
  no switchport
  ip address 10.1.11.1/31
  ip ospf network point-to-point
  ip router ospf 1 area 0.0.0.0
  ip ospf bfd
  mtu 9216
  no shutdown
router ospf 1
  router-id 10.1.1.1
\end{lstlisting}
\caption{BFD and OSPF configuration on Leaf1-1}
\label{fig:scenario1-bfd-ospf}
\end{figure}

\paragraph{Expected Behavior:}~\\
The adjacency between OSPF neighbors moves to FULL in all Spine-Leaf links. All loopback addresses, whether Loopback 0 BGP IDs or Loopback 1 VTEP addresses, appear in the OSPF routing table in every device. In every Leaf device, the OSPF routing table reflects dual equal-cost routes for each remote loopback address using each Spine device.

\paragraph{Results and Analysis}~\\
\begin{itemize}
    \item []Spine 1 OSPF neighbor table:
    \begin{figure}[!htbp]
        \centering
        \OspfSpine
        \caption{OSPF neighbor table on Spine 1 showing FULL adjacencies}
        \label{fig:ospf-spine}
    \end{figure}
    \item []Loopback reachability verification:
    \begin{figure}[!htbp]
        \centering
        \loopbackreachability
        \caption{Loopback reachability ping output}
        \label{fig:loopback-reachability}
    \end{figure}
    \newpage
    \item []Leaf 1 — OSPF routing table excerpt:
    \begin{figure}[!htbp]
        \centering
        \OspfRouting
        \caption{OSPF routing table on Leaf 1 showing ECMP paths over both spines}
        \label{fig:ospf-routing-ecmp}
    \end{figure}
\end{itemize}
The full state in the OSPF adjacency table means that bidirectional Hello exchange is successful on all the links of the spine-leaf and spine-border leaf connections. The routing table on Leaf 1 contains two equal cost paths for each loopback destination: a path over Spine 1 and a path over Spine 2.
\paragraph{Conclusion:}~\\
Full convergence has been achieved between the two sites. All the OSPF adjacencies are in FULL state, VTEP loopbacks are accessible from each leaf node via two equal cost paths, and ECMP is working on all the Leaf nodes. The problem of STP-caused bandwidth waste is solved all the links in the fabric are active.

\subsection{Scenario 2: BGP EVPN Peering and L2 Reachability}
\paragraph{Objective:}~\\
Ensure that there are iBGP EVPN sessions between all the leaf switches using the spine route reflectors, ensure that VXLAN VNIs are properly associated with the NVE interfaces, and ensure that two end hosts belonging to the same VNI connected to two different leaf switches can communicate without any data flooding.

\paragraph{Topology Context:}~\\
There are two virtual machines connected to two leaf switches at Site 1 VM-T1-1 connected to leaf 1 and VM-T1-2 connected to leaf 2. Both the virtual machines belong to VNI 10010. The spine switches are the iBGP Route Reflectors and reflect the EVPN routes to all the leaf switches acting as RR Clients.

\paragraph{Configuration Applied:}~\\
\begin{figure}[!htbp]
\centering
\begin{lstlisting}[
    language=bash,
    style=nexus,
]
router bgp 65100
  router-id 10.1.0.1
  address-family l2vpn evpn
    retain route-target all
  neighbor 10.1.1.0
    remote-as 65100
    update-source loopback0
    address-family l2vpn evpn
      send-community extended
      route-reflector-client
  neighbor 10.1.9.0
    remote-as 65100
    update-source loopback0
    address-family l2vpn evpn
      send-community extended
      route-reflector-client
\end{lstlisting}
\caption{BGP Route Reflector configuration on Spine 1}
\label{fig:scenario2-bgp-rr}
\end{figure}

\begin{figure}[!htbp]
\centering
\begin{lstlisting}[
    language=bash,
    style=nexus,
]
vlan 10
  vn-segment 10010

interface nve1
  no shutdown
  host-reachability protocol bgp
  source-interface loopback1
  member vni 10010
    ingress-replication protocol bgp

router bgp 65100
  router-id 10.1.1.1
  address-family l2vpn evpn
  neighbor 10.1.0.1
    remote-as 65100
    update-source loopback0
    address-family l2vpn evpn
      send-community extended
  neighbor 10.1.0.2
    remote-as 65100
    update-source loopback0
    address-family l2vpn evpn
      send-community extended

evpn
  vni 10010 l2
    rd auto
    route-target import auto
    route-target export auto
\end{lstlisting}
\caption{BGP EVPN and NVE configuration on Leaf 1}
\label{fig:scenario2-bgp-evpn-nve}
\end{figure}

\paragraph{Expected Behavior:}~\\
The iBGP EVPN sessions are up and are in the Established state for all the leaf switches. There is an NVE Peer table displaying all the remote VTEPs learned from Type 3 Inclusive Multicast Routes. When VM-T1-1 pings VM-T1-2, there is no flooding because the source leaf switch knows the MAC address of VM-T1-2 through the BGP EVPN Type 2 route advertisement.

\paragraph{Results and Analysis}~\\
\begin{itemize}
    \item []Leaf 1 BGP EVPN session summary:
    \begin{figure}[!htbp]
        \centering
        \EvpnLeaf
        \caption{BGP EVPN session summary on Leaf 1}
        \label{fig:evpn-leaf-summary}
    \end{figure}
    \item []Leaf 1 NVE peer table:
    \begin{figure}[!htbp]
        \centering
        \nvepeers
        \caption{NVE peer table showing remote VTEPs learned via EVPN}
        \label{fig:nve-peers}
    \end{figure}
    \item []Leaf 1 EVPN route table:
    \begin{figure}[!htbp]
        \centering
        \EVPNroute
        \caption{EVPN Type 2 route advertisement for VM-T1-2}
        \label{fig:evpn-route}
    \end{figure}
    \item []VM-T1-1 to VM-T1-2 ping result:
    \begin{figure}[!htbp]
        \centering
        \crosssite
        \caption{Successful ping between VMs across different leaf switches}
        \label{fig:cross-site-ping}
    \end{figure}
\end{itemize}
Verification confirms that Leaf 1 maintains stable iBGP sessions with both Spine Route Reflectors, successfully exchanging 28 prefixes to ensure correct EVPN route reflection across the fabric. The NVE peer table shows all remote VTEPs were discovered purely via control-plane (CP) Type 3 advertisements, eliminating data-plane discovery floods. Additionally, the EVPN table displays the proactive exchange of VM-T1-2's MAC/IP binding from Leaf 2 before traffic began, enabling Leaf 1 to route the ping directly to the destination VTEP without any initial flooding.
\paragraph{Conclusion:}~\\
Proper L2 reachability across the VXLAN fabric is established without any data plane flooding. The problem of L2 reachability instability.

\subsection{Scenario 3: ARP Suppression Verification}
\paragraph{Objective:}~\\
Prove that ARP requests made by locally connected endpoints are caught by the local VTEP and are resolved using the ARP cache populated by EVPN without flooding into the fabric as BUM traffic.

\paragraph{Topology Context:}~\\
VM-T1-1 residing on Leaf 1 makes an ARP request for the IP address of VM-T1-2 that resides on Leaf 2. Both VM-T1-1 and VM-T1-2 reside within VNI 10010. ARP suppression is configured on the NVE interface of Leaf 1.

\paragraph{Configuration Applied:}~\\
\begin{figure}[!htbp]
\centering
\begin{lstlisting}[
    language=bash,
    style=nexus,
]
interface nve1
  member vni 10010
    suppress-arp
\end{lstlisting}
\caption{ARP suppression configuration on Leaf 1 NVE interface}
\label{fig:scenario3-arp-suppression}
\end{figure}

\paragraph{Expected Behavior:}~\\
VM-T1-1 sends an ARP request for 192.168.10.2, which Leaf 1 intercepts and checks against its local suppression cache populated by Leaf 2's Type 2 EVPN routes. Leaf 1 then replies directly with VM-T1-2's MAC address, preventing the ARP broadcast from flooding any other VTEPs in the fabric.
\paragraph{Results and Analysis}~\\
\begin{itemize}
    \item []Leaf 1 ARP suppression cache:
    \begin{figure}[!htbp]
        \centering
        \ARPsuppression
        \caption{ARP suppression cache on Leaf 1 containing remote endpoints}
        \label{fig:arp-suppression-cache}
    \end{figure}
    \newpage
    \item []Leaf 1 NVE BUM traffic counters:
    \begin{figure}[!htbp]
        \centering
        \BUMEVPNTRAFFIC
        \caption{NVE BUM traffic counters showing zero BUM packets sent}
        \label{fig:bum-counters}
    \end{figure}
\end{itemize}
Leaf 1's ARP suppression cache maps all remote VNI 10010 endpoints to their respective VTEP tunnels via BGP EVPN Type 2 advertisements. Generating zero sent BUM packets proves that Leaf 1 resolved all local ARP requests internally, completely eliminating fabric-wide flooding. The 14 received BUM packets represent only the initial Type 3 VTEP membership exchanges during fabric bootup. This performance validates that the control plane successfully eliminates broadcast flooding scalability overhead.
\paragraph{Conclusion:}~\\
ARP suppression is working flawlessly. ARP requests are answered locally at the ingress VTEP and do not require broadcast flooding across the entire fabric.

\subsection{Scenario 4: L2 Extension and VM Mobility Across Sites}
\paragraph{Objective:}~\\
Ensure that the same VNI can be reached through the DCI link from Site 1 and Site 2, and confirm the BGP EVPN MAC Mobility feature used for detecting and resolving any VM migration.

\paragraph{Topology Context:}~\\
The VM-T1-1 and VM-T1-3 have connections to Sites 1 Leaf 1 and Sites 2 Leaf 1, respectively. Both use the VNI 10010. The DCI link spans the Border Gateway Leaf devices, FortiGate firewalls, and CSR1000v WAN router. The NDFC Multi-Site Domain concept applies the VNI 10010 in both sites.

\paragraph{Configuration Applied:}~\\
\begin{figure}[!htbp]
\centering
\begin{lstlisting}[
    language=bash,
    style=nexus,
]
evpn multisite border-gateway 100
  delay-restore time 300

interface nve1
  multisite border-gateway interface loopback1

router bgp 65100
  neighbor 10.3.9.1
    remote-as 65101
    address-family l2vpn evpn
      send-community extended
      rewrite-evpn-rt-asn
\end{lstlisting}
\caption{BGP EVPN multi-site border gateway configuration}
\label{fig:scenario4-bgp-evpn-ms}
\end{figure}

\paragraph{Expected Behavior:}~\\
VM-T1-1 in Site 1 connects to VM-T1-3 in Site 2 using VXLAN tunnel created via the DCI link. The BGP table of both sites shows Type 2 and Type 3 routes for remote VTEPs. For the purpose of simulating VM-T1-1 migrating from Leaf 1 to site 2 Leaf 1, which would constitute a cross-site vMotion, site 2 Leaf 1 generates a new Type 2 route using the same MAC address with MAC Mobility Sequence Number 1. All remote VTEPs then make necessary changes in their forwarding tables such that the connectivity for VM-T1-1 during migration remains intact without any manual intervention.

In practice, this migration process will be orchestrated by vCenter via a vMotion procedure. This involves moving VM's memory state, storage, and network identity from one ESXi host to another. It should be noted that the process of EVPN MAC Mobility is identical no matter whether migration is manually initiated or via vCenter as the new Type 2 announcement received by the remote VTEPs is treated identically.

\paragraph{Results and Analysis}~\\
\begin{itemize}
    \item []Leaf 1 site 1 EVPN route before migration:
    \begin{figure}[!htbp]
        \centering
        \beformigration
        \caption{EVPN Type 2 route before VM migration}
        \label{fig:evpn-before-migration}
    \end{figure}
    \item []Leaf 1 site 2 EVPN route after migration:
    \begin{figure}[!htbp]
        \centering
        \aftermigration
        \caption{EVPN Type 2 route after VM migration}
        \label{fig:evpn-after-migration}
    \end{figure}
    \item []Leaf 1 forwarding database update:
    \begin{figure}[!htbp]
        \centering
        \forwardingdatabase
        \caption{MAC forwarding table updated to remote VTEP after migration}
        \label{fig:forwarding-database}
    \end{figure}
    \item []Cross-site ping verification:
    \begin{figure}[!htbp]
        \centering
        \crosssite
        \caption{Successful cross-site ping after VM migration}
        \label{fig:cross-site-ping-after-vmotion}
    \end{figure}
\end{itemize}
In addition, based on the route tables generated in the EVPN environment, it can be observed that the VNI 10010 is successfully extended across both sites through the DCI link and the NDFC MSD topology. Increment in the MAC Mobility Sequence Number value from 0 to 1 within the Type 2 route announcement proves that the BGP EVPN control plane has successfully detected the VM migration event and the reachability information is updated in all the participants in the network fabric. The MAC table entries of the migrated VM on Leaf 1 shows the forwarding of the traffic toward Site 2 VTEP 10.4.1.1, indicating that there is automatic updating of the forwarding database within the local leaf device. Ping from Leaf 1 to site 2 Leaf 2-1 has been successful, confirming that the DCI link is active, and geographical L2 extension is working.

\paragraph{Conclusion:}~\\
Geographical L2 extension has been achieved using the above technique. The MAC Mobility approach automatically resolves VM migration events using the BGP EVPN control plane. Both constraints discussed in Chapter 1 regarding geographical L2 extension and VM mobility are overcome by implementing the same.

\subsection{Scenario 5: Inter-Site L3 Routing}
\paragraph{Objective:}~\\
Confirm that IP prefix reachability is set up between various subnets from Site 1 and Site 2 using BGP EVPN Type 5 IP Prefix routes and that inter-site routed traffic is taken via the optimal route through Border Gateway Leaf devices.

\paragraph{Topology Context:}~\\
The VM-T1-1 device on Site 1 (IP 192.168.10.1 in VRF-TELECOM) has to communicate with the VM-T2-4 device on Site 2 (IP 192.168.20.4 in VRF-WEBHOST subnet 192.168.20.0/24). The communication route should go through the DCI link through the Border Gateway Leaf and the CSR1000v WAN router where the Type 5 routes provide prefix reachability between the two sites.

\paragraph{Configuration Applied:}~\\
\begin{figure}[!htbp]
\centering
\begin{lstlisting}[
    language=bash,
    style=nexus,
]
router bgp 65100
  vrf VRF-TELECOM
    address-family ipv4 unicast
  vrf VRF-WEBHOST
    address-family ipv4 unicast
\end{lstlisting}
\caption{BGP VRF configuration on Border Gateway Leaf for inter-site routing}
\label{fig:scenario5-bgp-vrf}
\end{figure}

\paragraph{Expected Behavior:}~\\
Type 5 IP Prefix routes for site subnets are synchronized across the BGP EVPN tables of both locations, allowing each VRF routing table to learn remote prefixes via the Border Gateway VTEPs. Inter-site traffic is forwarded using Symmetric IRB. The ingress Leaf routes the packet into the tenant L3VNI, and the Border Gateway handles the cross-site encapsulation and decapsulation over the DCI link.
\paragraph{Results and Analysis}~\\
\begin{itemize}
    \item []Leaf 1 Type 5 routes in BGP EVPN table:
    \begin{figure}[!htbp]
        \centering
        \Typeroutes
        \caption{BGP EVPN Type 5 IP prefix routes on Leaf 1}
        \label{fig:type5-routes}
    \end{figure}
    \item []Leaf 1 VRF routing table:
    \begin{figure}[!htbp]
        \centering
        \VRFrouting
        \caption{VRF routing table showing remote site prefixes via EVPN}
        \label{fig:vrf-routing}
    \end{figure}
    \item []Inter-site traceroute:
    \begin{figure}[!htbp]
        \centering
        \intersite
        \caption{Traceroute output confirming optimal DCI path}
        \label{fig:intersite-traceroute}
    \end{figure}
\end{itemize}
The BGP EVPN table on Leaf 1 is able to populate Type 5 IP Prefix routes for remote site subnets that belong to the Border Gateway VTEP (10.2.9.1), indicating successful DCI prefix distribution. Moreover, the VRF routing table is able to list the remote prefixes, indicating that the Symmetric IRB configuration is able to route inter-site traffic through the transit L3VNI successfully. Traceroute outputs indicate an optimized data plane path comprising the ingress Leaf Anycast Gateway, Border Leaf DCI hand-off, CSR1000v WAN Transit, and remote Border Leaf, ensuring five hops.
\paragraph{Conclusion:}~\\
The Inter-Site L3 Routing with BGP EVPN Type 5 routes functionality is fully working. Remote prefixes are accessible via the optimal DCI link with minimal predictable and consistent latency. In addition, the inter-site routing inefficiency is fixed.

\subsection{Scenario 6: VRF Isolation and Inter-VRF Routing}
\paragraph{Objective:}~\\
Proof that traffic between tenants on the same site is fully isolated and separated at both L2 and L3 layers by utilizing the VRF and VNI design; proof of possibility to perform controlled communication between tenants using a specific Route Target import configuration.

\paragraph{Topology Context:}~\\
In total, there are three VMs attached to Leaf 1 on Site 1; one for each tenant, each residing in its own VRF: VM-T1-1 belongs to VRF-TELECOM (192.168.10.1), VM-T2-1 belongs to VRF-WEBHOST (192.168.20.1), and finally, VM-T3-1 belongs to VRF-CLIENT (192.168.30.1). The process is implemented in two steps: firstly, proving isolation, then controlling communication through a deliberate route leak.

\paragraph{Configuration Applied:}~\\
\begin{figure}[!htbp]
\centering
\begin{lstlisting}[
    language=bash,
    style=nexus,
]
vrf context VRF-TELECOM
  vni 50010
  rd auto
  address-family ipv4 unicast
    route-target import 65100:50010
    route-target export 65100:50010
    route-target import evpn 65100:50010
    route-target export evpn 65100:50010

vrf context VRF-WEBHOST
  vni 50020
  rd auto
  address-family ipv4 unicast
    route-target import 65100:50020
    route-target export 65100:50020
    route-target import evpn 65100:50020
    route-target export evpn 65100:50020
\end{lstlisting}
\caption{VRF configuration on Leaf switches}
\label{fig:scenario6-vrf-config}
\end{figure}

\begin{figure}[!htbp]
\centering
\begin{lstlisting}[
    language=bash,
    style=nexus,
]
vrf context VRF-TELECOM
  address-family ipv4 unicast
    route-target import 65100:50020
\end{lstlisting}
\caption{Route Target import for controlled cross-tenant communication}
\label{fig:scenario6-route-target-import}
\end{figure}

\paragraph{Expected Behavior:}~\\
In the first phase, the ping operation between VM-T1-1 and VM-T2-1 fails because of the complete absence of cross-tenant routing in the VRFs and VLAN separation. On the other hand, in phase 2, after the Route Target import from VRF-WEBHOST into VRF-TELECOM, the ping works successfully, and the WEBHOST subnet is reachable from the TELECOM tenant via the route leak.

\paragraph{Results and Analysis}~\\
\begin{itemize}
    \item []Phase 1 Cross-tenant ping failure:
    \newpage
    \begin{figure}[!htbp]
        \centering
        \corsstenatfa
        \caption{Phase 1: Cross-tenant ping failure (100\% packet loss)}
        \label{fig:phase1-ping-fail}
    \end{figure}
    \item []Leaf 1 VRF routing table Phase 1:
    \begin{figure}[!htbp]
        \centering
        \pahesI
        \caption{VRF routing table Phase 1 – no cross-tenant routes}
        \label{fig:phase1-vrf-table}
    \end{figure}
    \item []Phase 2 Cross-tenant ping success after route leak:
    \begin{figure}[!htbp]
        \centering
        \corsstenat
        \caption{Phase 2: Cross-tenant ping success after route leak}
        \label{fig:phase2-ping-success}
    \end{figure}
    \item []Leaf 1 VRF routing table Phase 2:
    \begin{figure}[!htbp]
        \centering
        \phaseII
        \caption{VRF routing table Phase 2 – WEBHOST subnet injected}
        \label{fig:phase2-vrf-table}
    \end{figure}
    \item []VRF-CLIENT isolation confirmed:
    \begin{figure}[!htbp]
        \centering
        \VRFCLIENT
        \caption{VRF-CLIENT remains isolated from TELECOM}
        \label{fig:vrf-client-isolation}
    \end{figure}
\end{itemize}
Phase 1 findings confirm total tenant isolation: each VRF routing table maintains only its local subnet, and cross-tenant pings to VRF-WEBHOST fail with 100\% packet loss. This proves that the VNI/VRF combination enforces strict isolation at both Layer 2 and Layer 3, surpassing the 4,096 VLAN limit to deliver the granular multi-tenancy required for Tunisie Telecom's diverse clientele. Phase 2 results demonstrate that this isolation is highly controllable; configuring a single Route Target import safely exposed the WEBHOST subnet to TELECOM while remaining hidden from CLIENT. Ultimately, leveraging the 24-bit VNI namespace provides over 16 million logical segments, completely solving the infrastructure's scalability challenges.
\paragraph{Conclusion:}~\\
The isolation of tenants in question is achieved at both Layer 2 and Layer 3 using the VNI/VRF implementation. Control of cross-tenant traffic is also possible with Route Targets' configuration.

\subsection{Scenario 7: NDFC Fabric Onboarding and Compliance Verification}
\paragraph{Objective:}~\\
Demonstrate that the NDFC offers full automation for BGP EVPN-VXLAN fabric provisioning from discovery through underlay and overlay deployment, and verify that the compliance engine detects and logs configuration drift caused by in-band management.

\paragraph{Topology Context:}~\\
NDFC manages all the Site 1 devices including the Spine devices, the Leaf devices, and the Border Gateway Leaf devices in the in-band management plane. This is due to the fact that all the devices can be accessed from the NDFC via the Loopback 100 interface that is associated with VRF-MGMT, where all the routes are advertised using the OSPF underlay throughout the fabric. The NDFC device connects to the network infrastructure via the Leaf 1-X node using VLAN 99, while the management subnet 192.168.100.0/24 is expanded to Site 2 via Type 5 route advertisements using VRF-MGMT.

\paragraph{Configuration Applied:}~\\
Fabric provisioning happens using the NDFC Easy Fabric process and not using any manual configuration for each device individually. The NDFC creates the entire underlay and overlay configurations using the “Recalculate \& Deploy” feature.

\begin{figure}[!htbp]
\centering
\begin{lstlisting}[
    language=bash,
    style=nexus,
]
Fabric name:       TT-SITE1-FABRIC
BGP AS:            65100
OSPF process:      1
Underlay routing:  OSPF
Overlay protocol:  BGP EVPN
iBGP peering:      Spine as RR
Replication mode:  Ingress
Anycast GW MAC:    0200.5e00.01ff
BFD:               Enabled (300ms x 3)
Management:        In-band VRF-MGMT
                   VLAN 99
                   192.168.100.0/24
\end{lstlisting}
\caption{NDFC Easy Fabric settings applied}
\label{fig:scenario7-easy-fabric}
\end{figure}

\paragraph{Results and Analysis}~\\
\begin{itemize}
    \item []NDFC fabric topology view:
    \begin{figure}[!htbp]
        \centering
        \fabric
        \caption{NDFC fabric topology view showing all discovered devices}
        \label{fig:ndfc-topology}
    \end{figure}
    \item []NDFC fabric dashboard:
    \begin{figure}[!htbp]
        \centering
        \fabricdash
        \caption{NDFC fabric dashboard with In-Sync status}
        \label{fig:ndfc-dashboard}
    \end{figure}
\end{itemize}
Indeed, the NDFC fabric topology clearly indicates the discovery and role configuration of all nine Cisco Nexus 9000 devices: two Spine routers, seven Leafs, and two border gateways. The consistency in the In-Sync status displayed in the dashboard clearly shows that there is a 100\% match between the currently running configurations and the ones stored by NDFC.

The Recalculate and Deploy procedure has fully automated the deployment process for the OSPF underlay network, loopback interfaces, iBGP EVPN peers, NVE interfaces, and VNI bindings without any manual commands through the CLI. Additionally, the compliance engine was able to detect an intentional modification performed outside of NDFC on Leafs 1-3, confirming its capability to detect changes made on the devices.

\paragraph{Conclusion:}~\\
NDFC is a complete automation solution for the entire fabric lifecycle. It continuously enforces the intent and detects any deviation in the system.

\subsection{Scenario 8: Spine Failure and BFD Reconvergence}
\paragraph{Objective:}~\\
Make sure that the fabric has detected the failure of one of the Spine nodes using BFD inside the detection window configured, that the fabric uses OSPF to reconverge through the other Spine node, and that the end-to-end traffic is recovered inside a sub-second window measuring the time saved compared to BGP hold timer mechanism.

\paragraph{Topology Context:}~\\
An iperf3 traffic flow operates continually from VM-T1-1 running on Leaf 1 and VM-T1-2 running on Leaf 2 at Site 1. Both of the Leaf nodes have established OSPF neighbor relationships and ECMP routes through both of the Spines 1 and 2. BFD functionality has been enabled on all OSPF sessions, and detection timer is set to be 300 milliseconds while a multiplier equals to 3.

\paragraph{Configuration Applied:}~\\
(No additional configuration beyond the global BFD settings shown in Scenario 1.)

\paragraph{Expected Behavior:}~\\
While the iperf3 traffic is flowing continuously, administrative shutdown of the first of the Spines occurs. After detecting the absence of any adjacent routers using BFD (within 900 milliseconds), the information is passed to the OSPF protocol that withdraws all of the routes advertised by Spine 1 and establishes the rest of routes via Spine 2.

\paragraph{Results and Analysis}~\\
\begin{itemize}
    \item []BFD neighbor state and OSPF neighbors before failure:
    \begin{figure}[!htbp]
        \centering
        \bfdsucc
        \caption{BFD neighbor state before spine failure}
        \label{fig:bfd-before}
    \end{figure}
    \begin{figure}[!htbp]
        \centering
        \Ospfleaf
        \caption{OSPF neighbor table before spine failure}
        \label{fig:ospf-before}
    \end{figure}
    \item []BFD neighbor state and OSPF neighbors after failure:
    \begin{figure}[!htbp]
        \centering
        \bfdfail
        \caption{BFD neighbor state after spine failure}
        \label{fig:bfd-after}
    \end{figure}
    \begin{figure}[!htbp]
        \centering
        \Ospfleaffailer
        \caption{OSPF neighbor table after reconvergence}
        \label{fig:ospf-after}
    \end{figure}
    \newpage
    \item[]iperf3 traffic recovery:
    \begin{figure}[!htbp]
        \centering
        \Ipsec
        \caption{iperf3 traffic recovery within one second after spine failure}
        \label{fig:iperf-recovery}
    \end{figure}
\end{itemize}
The BFD neighbor table reveals that the session to Spine 1 terminated immediately due to failure, while the session to Spine 2 continued to be in the FULL state, demonstrating that BFD works separately for each adjacency. Following reconvergence, the OSPF neighbor table and routing table reveal that all routes have been removed from Spine 1 and switched to Spine 2, with the ECMP path removed.

iperf3 measurements reveal that for just one second the BFD/OSPF reconvergence caused zero throughput, after which it recovered immediately. This quick sub-second convergence (limited by 900ms BFD detection period) is an enormous enhancement compared to the standard BGP 90 second Hold timer. This addresses the availability issues and single point of failure risks.

\paragraph{Conclusion:}~\\
Spine failure detection is achieved in 900ms, resulting in OSPF reconvergence immediately. Traffic converges within one measurement second. This leads to the detection time improvement by 99\% compared to BGP. This resolves the availability issues.

\section{Conclusion}
This chapter has tested the whole BGP EVPN-VXLAN IP Fabric architectural design in terms of eight distinct test cases, with each test case focused on testing a particular feature of the architecture to solve a pain point stated in Chapter 1.

In summary, all the findings from the tests show that the architecture under study meets every requirement set by this project since the beginning. This architecture has enabled Tunisie Telecom to implement an efficient and scalable DC network that is not limited by any shortcomings of the traditional three-tier architecture and can handle AI-powered computing workloads in the future.